\section{Derivation of Expected Precision Goal Stop Iteration}\label{app:expected_stop}

The precision-based stopping algorithms (PitG and ePitG) halt when the HDI width
(see Appendix~\ref{app:hdi})
falls below the precision goal $\omega_{\rm goal}$. Because the posterior HDI width
shrinks as $N^{-1/2}$, one can predict in advance the iteration at which this
threshold will be crossed. The derivation below uses a normal approximation to the
posterior and is therefore intended as an analytical guide rather than an exact
finite-sample identity.
% Throughout, we assume unimodal, approximately symmetric posteriors, for which
% HDI width is a natural and interpretable precision measure.
% For multimodal posteriors, HDI width becomes ambiguous (it is unclear which mode
% to centre on), and posterior entropy may be a more appropriate precision criterion
% \citep{kruschke2015doing, chaloner1995, sebastiani2000}; such settings fall outside
% the scope of the present work.

\subsection*{General CLT Form}

Assume an estimator $\hat\theta_N$ with asymptotic distribution
\begin{equation}\label{eq:clt_general}
\hat\theta_N \approx \mathcal{N}\!\left(\theta,\,\frac{V(\theta)}{N}\right),
\end{equation}
where $V(\theta)$ is the per-observation variance contribution. For a symmetric
central interval with critical value $z_{*}$ (e.g., $z_{*}=1.96$ for 95\%), the
interval width is approximately
\begin{equation}\label{eq:hdi_width_approx}
w \approx 2 z_{*}\sqrt{\frac{V(\theta)}{N}}.
\end{equation}
Setting $w = \omega_{\rm goal}$ and solving for $N$ gives the general expected stop
iteration:
\begin{equation}\label{eq:expected_stop_iteration_clt}
N_{\rm goal} \approx \frac{4 z_{*}^{2}}{\omega_{\rm goal}^{2}}\,V(\theta).
\end{equation}

\subsection*{Binomial Case}

For Bernoulli outcomes, $V(\theta)=\theta(1-\theta)$, so
Equation~\ref{eq:expected_stop_iteration_clt} becomes
\begin{equation}\label{eq:expected_stop_iteration_binomial}
N_{\rm goal} \approx \frac{4 z_{*}^{2}}{\omega_{\rm goal}^{2}}\,\theta(1-\theta).
\end{equation}
This is maximised at $\theta=0.5$, confirming that the fair coin case requires the
largest sample to achieve a given precision goal, and decreases monotonically as
$\theta$ moves towards $0$ or $1$.

\subsection*{Continuous Case (Mean Estimation)}

For continuous outcomes when the estimand is a mean $\mu$ with variance $\sigma^{2}$,
$V(\mu)=\sigma^{2}$ (or a plug-in estimate $\hat\sigma^{2}$), giving
\begin{equation}\label{eq:expected_stop_iteration_mean}
N_{\rm goal} \approx \frac{4 z_{*}^{2}\sigma^{2}}{\omega_{\rm goal}^{2}}.
\end{equation}

\subsection*{Indexing Convention}

If updates are indexed by $n=0,1,2,\dots$ while sample size is $N=n+1$, rewriting
Equation~\ref{eq:expected_stop_iteration_clt} in terms of $n$ gives
\begin{equation}\label{eq:expected_stop_index}
n_{\rm goal} \approx \frac{4 z_{*}^{2}}{\omega_{\rm goal}^{2}}\,V(\theta) - 1.
\end{equation}
This is an indexing convention only; asymptotically, the distinction is negligible.

Figure~\ref{fig:min_sample_by_goal} in Section~\ref{sec:expected_stop_iteration} of the
main text illustrates $N_{\rm goal}$ as a function of $\theta_{\rm true}$ for both the
Bernoulli and continuous cases across a range of precision goals.

\subsection*{Suggested References}

The approximation above rests on standard asymptotic results. Relevant sources include:
\cite{kruschke2015doing} (Ch.\ 13, PitG context);
van der Vaart (1998), \textit{Asymptotic Statistics} (CLT and asymptotic normality);
Casella \& Berger (2002), \textit{Statistical Inference} (CI-width sample-size heuristics);
Brown, Cai \& DasGupta (2001), \textit{Statistical Science} (binomial proportion intervals);
Chow, Shao \& Wang (2008), \textit{Sample Size Calculations in Clinical Research}
(practical width-based formulas).
